% ============================================================================
% preamble.tex — 《如何用 AI Agent 编写自己的操作系统》
% ============================================================================

% --- 中文支持 (XeLaTeX) ---
\usepackage{ctex}
\usepackage{fontspec}

% --- 页面布局 ---
\usepackage[
  a4paper,
  top=2.5cm, bottom=2.5cm,
  left=2.5cm, right=2.5cm
]{geometry}

% --- 数学 ---
\usepackage{amsmath, amssymb}

% --- 图表 ---
\usepackage{graphicx}
\usepackage{float}
\usepackage{booktabs}
\usepackage{longtable}
\usepackage{tabularx}

% --- 绘图 ---
\usepackage{tikz}
\usetikzlibrary{
  positioning, calc, arrows.meta,
  decorations.pathreplacing, fit,
  shapes.multipart, shapes.geometric,
  patterns, backgrounds,
}

% Agent 工作流图专用样式
\tikzset{
  agentbox/.style={
    draw, rounded corners=4pt, minimum width=2.8cm, minimum height=1cm,
    font=\small\sffamily, align=center, thick,
  },
  humanbox/.style={agentbox, fill=blue!10, draw=blue!50!black},
  aibox/.style={agentbox, fill=orange!10, draw=orange!50!black},
  artifactbox/.style={
    draw, dashed, rounded corners=2pt, minimum width=2.2cm, minimum height=0.7cm,
    font=\small\ttfamily, fill=gray!5,
  },
  flowedge/.style={-{Stealth[length=5pt]}, thick},
  parallededge/.style={-{Stealth[length=5pt]}, thick, dashed, blue!60},
}

% --- 颜色与代码高亮 ---
\usepackage{xcolor}
\usepackage{listings}

% 代码样式：C 语言
\lstdefinestyle{cstyle}{
  language=C,
  basicstyle=\ttfamily\small,
  keywordstyle=\color{blue}\bfseries,
  commentstyle=\color{gray}\itshape,
  stringstyle=\color{red!70!black},
  numbers=left,
  numberstyle=\scriptsize\color{black!40},
  numbersep=10pt,
  frame=l,
  framesep=6pt,
  framerule=1.5pt,
  rulecolor=\color{black!30},
  breaklines=true,
  tabsize=4,
  showstringspaces=false,
  columns=flexible,
}

% 代码样式：YAML / Markdown（agent 配置文件）
\lstdefinestyle{yamlstyle}{
  basicstyle=\ttfamily\small,
  keywordstyle=\color{blue}\bfseries,
  commentstyle=\color{gray}\itshape,
  stringstyle=\color{red!70!black},
  numbers=left,
  numberstyle=\scriptsize\color{black!40},
  numbersep=10pt,
  frame=l,
  framesep=6pt,
  framerule=1.5pt,
  rulecolor=\color{black!30},
  breaklines=true,
  tabsize=2,
  showstringspaces=false,
  columns=flexible,
  morekeywords={description, tools, name, model, agents, mode},
}

% 代码样式：x86 汇编
\lstdefinestyle{nasmstyle}{
  language=[x86masm]Assembler,
  basicstyle=\ttfamily\small,
  keywordstyle=\color{blue}\bfseries,
  commentstyle=\color{gray}\itshape,
  numbers=left,
  numberstyle=\scriptsize\color{black!40},
  numbersep=10pt,
  frame=l,
  framesep=6pt,
  framerule=1.5pt,
  rulecolor=\color{black!30},
  breaklines=true,
  tabsize=4,
  showstringspaces=false,
  columns=flexible,
}

% Shell
\lstdefinestyle{shellstyle}{
  language=bash,
  basicstyle=\ttfamily\small,
  keywordstyle=\color{blue},
  commentstyle=\color{gray}\itshape,
  frame=l,
  framesep=6pt,
  framerule=1.5pt,
  rulecolor=\color{black!30},
  breaklines=true,
  tabsize=4,
  showstringspaces=false,
  columns=flexible,
}

% Prompt / 对话记录样式
\lstdefinestyle{chatstyle}{
  basicstyle=\ttfamily\small,
  breaklines=true,
  tabsize=2,
  showstringspaces=false,
  columns=flexible,
  frame=l,
  framesep=6pt,
  framerule=1.5pt,
  rulecolor=\color{purple!40},
}

\lstset{style=cstyle}

% 终端样式
\definecolor{sol-base03}{HTML}{002B36}
\definecolor{sol-base02}{HTML}{073642}
\definecolor{sol-base0}{HTML}{B0BEC5}
\definecolor{sol-cream}{HTML}{EEE8D5}
\definecolor{sol-green}{HTML}{859900}
\definecolor{sol-cyan}{HTML}{2AA198}

\newcommand{\tin}[1]{{\color{sol-cream}\bfseries #1}}
\newcommand{\tps}[1]{{\color{sol-green}\bfseries #1}}
\newcommand{\tpk}[1]{{\color{sol-cyan}\bfseries #1}}

\lstdefinestyle{consolestyle}{
  basicstyle=\ttfamily\small\color{sol-base0},
  keywordstyle=\color{sol-base0},
  commentstyle=\color{sol-base0},
  stringstyle=\color{sol-base0},
  identifierstyle=\color{sol-base0},
  numbers=none,
  frame=none,
  breaklines=true,
  tabsize=4,
  showstringspaces=false,
  columns=flexible,
  escapeinside={«}{»},
}

% --- 超链接 ---
\usepackage{hyperref}
\hypersetup{
  colorlinks=true,
  linkcolor=blue!60!black,
  urlcolor=blue!80!black,
  citecolor=green!50!black,
  pdfauthor={SZY},
  pdftitle={如何用 AI Agent 编写自己的操作系统},
}

% --- 页眉页脚 ---
\usepackage{fancyhdr}
\pagestyle{fancy}
\fancyhf{}
\fancyhead[L]{\leftmark}
\fancyhead[R]{\thepage}
\renewcommand{\headrulewidth}{0.4pt}

% --- 自定义命令 ---
\newcommand{\reg}[1]{\texttt{#1}}
\newcommand{\hex}[1]{\texttt{0x#1}}
\newcommand{\term}[1]{\textbf{#1}}
\newcommand{\codefile}[1]{%
  \par\noindent{\small\ttfamily\color{blue!60!black}#1}\vspace{-0.5em}%
}
\newcommand{\agentfile}[1]{%
  \par\noindent{\small\ttfamily\color{orange!60!black}#1}\vspace{-0.5em}%
}

% Agent 角色标记
\newcommand{\agentarchitect}{\textsf{\textbf{@Architect}}}
\newcommand{\agentkernel}{\textsf{\textbf{@Kernel}}}
\newcommand{\agentauthor}{\textsf{\textbf{@Author}}}
\newcommand{\agentship}{\textsf{\textbf{/ship}}}

% --- tcolorbox ---
\usepackage{tcolorbox}
\tcbuselibrary{breakable, skins, listings}

% 终端输出框
\tcbset{
  consolebox/.style={
    colback=sol-base03,
    colframe=sol-base02,
    coltitle=sol-cream,
    fonttitle=\bfseries\ttfamily\small,
    title={\textcolor{sol-green}{\$}~Terminal},
    listing only,
    listing options={style=consolestyle},
    breakable,
    sharp corners,
    boxrule=0.5pt,
    left=4pt, right=4pt, top=2pt, bottom=2pt,
  }
}

% 设计决策框
\newtcolorbox{designbox}[1][]{
  colback=blue!3, colframe=blue!40!black,
  fonttitle=\bfseries, title={设计决策}, breakable, #1
}

% 经验教训框
\newtcolorbox{lessonbox}[1][]{
  colback=green!3, colframe=green!40!black,
  fonttitle=\bfseries, title={经验教训}, breakable, #1
}

% 反模式 / 踩坑框
\newtcolorbox{pitfallbox}[1][]{
  colback=red!3, colframe=red!40!black,
  fonttitle=\bfseries, title={踩坑警告}, breakable, #1
}

% Prompt 技巧框
\newtcolorbox{promptbox}[1][]{
  colback=purple!3, colframe=purple!40!black,
  fonttitle=\bfseries, title={Prompt 技巧}, breakable, #1
}

% 对话记录框（人类 / AI 交替）
\newtcolorbox{humanmsg}[1][]{
  colback=blue!5, colframe=blue!30,
  fonttitle=\bfseries\small, title={🧑 Human},
  breakable, left=4pt, right=4pt, #1
}
\newtcolorbox{aimsg}[1][]{
  colback=orange!5, colframe=orange!30,
  fonttitle=\bfseries\small, title={🤖 Agent},
  breakable, left=4pt, right=4pt, #1
}

% --- 思考题与练习 ---
\newcounter{thinknum}[chapter]
\newcounter{exercisenum}[chapter]

\newtcolorbox{thinkbox}[1][]{
  colback=purple!3, colframe=purple!40!black,
  fonttitle=\bfseries,
  title={\stepcounter{thinknum}思考 \thechapter.\thethinknum},
  breakable, #1
}

\newtcolorbox{exercise}[1][]{
  colback=cyan!3, colframe=cyan!40!black,
  fonttitle=\bfseries,
  title={\stepcounter{exercisenum}练习 \thechapter.\theexercisenum},
  breakable, #1
}
